\section{总结与延伸}

\subsection{核心要点回顾}

这篇 HumanLayer 的文章系统性地论述了 coding agent 性能优化的方法论。核心论点可以浓缩为一句话：\textbf{``The model is probably fine. It's just a skill issue.''}---当你的 coding agent 表现不佳时，问题大概率出在 harness（配置与运行环境），而非模型本身。

文章介绍的六大 harness 配置杠杆可以总结如下：

\begin{table}[H]
\centering
\caption{六大 Harness 配置杠杆总览}
\begin{tabular}{@{}p{3.5cm}p{4cm}p{5.5cm}@{}}
\toprule
\textbf{配置杠杆} & \textbf{核心作用} & \textbf{关键原则} \\
\midrule
CLAUDE.md / AGENTS.md & 基础行为配置 & 精简、普适、避免自动生成 \\
MCP Servers & 能力扩展 & 按需启用、CLI 优先 \\
Skills & 渐进式知识披露 & 按需加载、注意安全 \\
Sub-Agents & Context 隔离 & 用于隔离而非角色分工 \\
Hooks & 确定性控制流 & 成功静默、失败详报 \\
Back-Pressure & 自动验证 & 验证输出必须精简 \\
\bottomrule
\end{tabular}
\end{table}

\subsection{贯穿全文的设计理念}

文章反复强调的几个设计理念值得深思：

\begin{enumerate}[leftmargin=1.5em]
    \item \textbf{Context Efficiency}：所有配置决策的终极评判标准是``这是否让 context window 中的信息更高效？''
    \item \textbf{Progressive Disclosure}：信息按需加载，避免预加载所有可能需要的内容
    \item \textbf{Silent Success, Verbose Failure}：成功不产生输出，失败才详细报告---这是验证机制的设计金律
    \item \textbf{Reactive, Not Proactive}：不要预防性优化，在真实失败中学习和改进
    \item \textbf{Instruction Budget}：将 context window 视为有限预算，每条指令都有机会成本
\end{enumerate}

\subsection{拓展阅读}

\begin{itemize}[leftmargin=1.5em]
    \item HumanLayer 关于 CLAUDE.md 最佳实践的博文
    \item Dex 的 12-factor agents 框架
    \item Viv 关于 harness engineering 的系列文章
    \item ETH Zurich 的 agentfile 有效性研究（138 个 agentfile 的实证分析）
    \item Chroma 的 context rot 研究
    \item Terminal Bench 2.0 benchmark
    \item Matt Pocock 关于 AGENTS.md 的文章
    \item OpenAI 关于 harness engineering 的博文
\end{itemize}

%% --- Fin del contenido --- %%

\end{document}
